\section{Evaluasi Akhir dan Integrasi Kompetensi}

Bab terakhir ini merupakan titik di mana seluruh kompetensi dari bab-bab sebelumnya diintegrasikan dalam proyek akhir yang di-deploy ke lingkungan produksi. Evaluasi akhir mengukur pencapaian seluruh Sub-CPMK: kemampuan membangun front-end (HTML/CSS/JS), back-end (PHP, MySQL), integrasi API, keamanan web, pengujian, dan deployment. Proyek akhir bukan sekadar tugas teknis --- ini adalah simulasi proses pengembangan web yang sesungguhnya \cite{mdnDeploy}.

\begin{konsep}
\textbf{Kompetensi Terintegrasi}: Proyek akhir harus menunjukkan kemampuan menyeluruh:
\begin{itemize}
  \item \textbf{Front-end}: layout responsive, form validation, Fetch API
  \item \textbf{Back-end}: PHP OOP, session, CRUD dengan prepared statement
  \item \textbf{Database}: desain ternormalisasi, relasi FK, backup
  \item \textbf{Keamanan}: validasi input, XSS/CSRF prevention, HTTPS
  \item \textbf{Deployment}: konfigurasi produksi, permission, monitoring
\end{itemize}
\end{konsep}

Refleksi kompetensi: review checklist di setiap bab; identifikasi area yang masih perlu dikembangkan; buat rencana belajar lanjutan. Kelulusan mata kuliah menandai awal perjalanan --- teknologi web terus berkembang, dan seorang developer harus terus belajar.

